\documentclass{beamer}
\usetheme{Berlin}
\usepackage[latin1]{inputenc}
\usepackage{graphicx}
\usepackage{alltt}
\usepackage[spanish]{babel}
\title{Android Malware}
\author{Francisco Ar\'evalo \\ \texttt{francisco.arevalo@continuum.cl}}
\date{\today}
\begin{document}
\begin{frame}
\titlepage
\begin{figure}[htb]
\centering
\includegraphics[scale=0.3]{isotipo_continuum}
\end{figure}
\end{frame}
\section{Amenazas}
\subsection{Amenazas (1/4)}
\begin{frame}[shrink=3]
\frametitle{Amenazas de aplicaciones}
\begin{itemize}
\item Malware\\
Software dise\~nado para efectuar operaciones no deseadas en un dispositivo. Puede generar cargos a la cuenta, enviar mensajes, o controlar el funcionamiento del sistema.
\item Spyware\\
Software que esp\'ia o recolecta datos privados sin consentimiento del usuario, incluidos e-mail, mensajes, libretas de direcciones e im\'agenes.
\item Privacy leaks\\
Aplicaciones que divulgan informaci\'on privada como la ubicaci\'on, historial o aplicaciones instaladas de manera no intencional.
\item Aplicaciones Vulnerabiles\\
Programas con fallas de dise\~no, que permiten acceso no autorizado, o generan problemas con la usabilidad del dispositivo donde se instalan.
\end{itemize}
\end{frame}
\subsection{Amenazas (2/4)}
\begin{frame}
\frametitle{Amenazas web}
\begin{itemize}
\item Phishing\\
P\'aginas o formularios que suplantan o enga\~nan al usuario para que entregue informaci\'on confidencial a un atacante.
\item Drive-by Downloads\\
Descargas que se inician autom\'aticamente al visitar un sitio comprometido.
\item Browser exploits\\
Aprovechan vulnerabilidades del explorador o sus plugins para instalar malware o ejecutar comandos en un dispositivo.
\end{itemize}
\end{frame}
\subsection{Amenazas (3/4)}
\begin{frame}
\frametitle{Amenazas de red}
\begin{itemize}
\item Network exploits\\
Amenazas que utilizan fallas en los protocolos de Bluetooth, Wi-Fi, SMS, o red celular para tomar control de un dispositivo.
\item Sniffing Wi-Fi\\
Captura de datos no ecriptados, en especial passwords y emails, al transmitirlos por una red insegura
\end{itemize}
\end{frame}
\subsection{Amenazas (4/4)}
\begin{frame}
\frametitle{Riesgos F\'isicos}
\begin{itemize}
\item P\'erdida del equipo\\
El riesgo m\'as com\'un, ultimamente se han reportado casos de chantaje o abuso de informaci\'on privada que se encuentra en un tel\'efono extraviado.
\end{itemize}
\end{frame}
\section{Propagaci\'on}
\subsection{Propagaci\'on (1/2)}
\begin{frame}[shrink=2]
\frametitle{?`Por qu\'e la gente instala malware?}
\begin{itemize}
\item Repackaging\\
	{\scriptsize El malware se disfraza dentro de una aplicaci\'on leg\'itima y vuelta a publicar con un nombre e \'icono similar. Ocurre a menudo con aplicaciones 'crackeadas'.}
\item Bait and switch\\
	{\scriptsize La aplicaci\'on promete funcionalidades que en realidad no tiene, o se publica bajo una descripci\'on falsa.}
\item Ataques de Update\\
	{\scriptsize El autor del malware publica una aplicaci\'on leg\'itima, pero al actualizarla, descarga una versi\'on modificada con c\'odigo malicioso.}
\end{itemize}
\end{frame}
\subsection{Propagaci\'on (2/2)}
\begin{frame}[shrink=2]
\frametitle{?`Por qu\'e la gente instala malware?}
\begin{itemize}
\item Shotgun distribution\\
	{\scriptsize El malware se publica bajo diversos nombres, en varios app stores, y por distintos autores. En el caso de DroidDream, se crearon m\'as de 60 aplicaciones infectadas.}
\item Malvertising\\
	{\scriptsize El desarrollador aprovecha los 'in-app advertisement' para dirigir al usuario hacia un sitio de descarga infectado. Estos avisos pueden aparecer incluso dentro de apicaciones leg\'itimas.}
\item Drive-by Downloads\\
	{\scriptsize Al visitar un sitio infectado, se inicia la descarga del malware autom\'aticamente. Muchas veces se simulan notificaciones de 'trusted download', o di\'alogos para enga\~nar al usuario y lograr que ejecute el archivo descargado.}
\end{itemize}
\end{frame}
\section{Consecuencias}
\subsection{Consecuencias (1/2)}
\begin{frame}
\frametitle{?`Cu\'ales son las consecuencias?}
\begin{itemize}
\item SMS Premium:\\
	{ \scriptsize GGTracker\footnote{Security Alert: Android Trojan GGTracker Charges Premium Rate SMS Messages \texttt{http://bit.ly/P3lYP6}}, RuFraud\footnote{Security Alert: Android Trojan GGTracker Charges Premium Rate SMS Messages \texttt{http://bit.ly/rTUQIu}}, aprovechan que las suscripciones a estos servicios son la forma m\'as simple de generar cobros directamente al usuario.}
\item Botnets:\\
	{ \scriptsize DroidDream\footnote{Encyclopedia entry: TrojanSpy:AndroidOS/DroidDream.A \texttt{http://bit.ly/MKOwbL}}, convierten al equipo en un 'zombie', reciben instrucciones remotas y exponen todas las funcionalidades del tel\'efono (realizar llamadas, eviar mensajes, abrir conecciones a internet).}
\end{itemize}
\end{frame}
\subsection{Consecuencias (2/2)}
\begin{frame}
\frametitle{?`Cu\'ales son las consecuencias?}
\begin{itemize}
\item Elevaci\'on de privilegios:\\
	{\scriptsize Comunes al intentar 'rootear', o acceder al usuario root en un equipo. Una vez que el exploit obtiene acceso root, puede cargar aplicaciones, o acceder a cualquier funcionalidad del tel\'efono.}
\item Spyware y aplicaciones esp\'ia personalizadas:\\
	{\scriptsize Pueden enfocarse a robar cualquier informaci\'on que pase por el tel\'efono, o personalizarse para obtener credenciales bancarias o de servicios privados.}
\end{itemize}
\end{frame}
\section{Conclusiones}
\begin{frame}
\frametitle{Conclusiones}
\begin{itemize}
\item S\'olo descargar aplicaciones de fuentes confiables
\item Revisar que los links apunten a la direcci\'on correcta antes de hacer click
\item El buen manejo de claves y cuentas de usuario es indispensable
\item Utilizar herramientas de seguridad para dispositivos m\'oviles (Lookout, AVG, Norton, Trend Micro) 
\item Estar atento a comportamientos extra\~nos en las aplicaciones, y cargos inusuales en la cuenta mensual
\item Mantener el firmware actualizado
\item Revisar los permisos que solicitan las aplicaciones antes de instalar
\end{itemize}
\end{frame}
\section{Bibliograf\'ia}
\begin{frame}
\frametitle{Links y contacto}
\begin{itemize}
\item Lookout Mobile Security \\ \texttt{https://www.mylookout.com/}
\item Android Genome Project \\ \texttt{http://lwn.net/Articles/498698/}
\item IEEE Symposium on Security \& Privacy \\ \texttt{http://www.ieee-security.org/TC/SP2012/}\\
\end{itemize}
\begin{figure}[htb]
\centering
\texttt{francisco.arevalo@continuum.cl}
\end{figure}
\end{frame}
\end{document}
